% arara: clean: {extensions: ['log','aux','synctex.gz','out']}
% arara: pdflatex: { shell: yes }
% arara: pdflatex: { shell: yes }
% arara: clean: {extensions: ['log','aux','synctex.gz','out']}
\documentclass[spanish,addpoints,answers,a4paper]{exam}
\usepackage[T1]{fontenc}
\usepackage[spanish]{babel}
\usepackage{libertine}
\usepackage[tmargin=2cm,rmargin=2cm,lmargin=2cm]{geometry}
\usepackage{amsmath,bm}
\usepackage[shortlabels]{enumitem}
\usepackage{siunitx}
\sisetup{
	group-four-digits = true,
	group-separator = {,}
}
\usepackage{multicol}

\usepackage{graphicx}
\graphicspath{{../../../images/}}
\pagestyle{plain}

\usepackage{minted}
\setminted[cpp]{tabsize=2}
\usepackage{hyperref}
\renewcommand{\solutiontitle}{\noindent\textbf{Solución}\par\noindent}
\renewcommand\listingscaption{Listado}
\newcommand{\unmarkedfntext}[1]{%
	\begingroup
	\renewcommand\thefootnote{}\footnote{#1}%
	\addtocounter{footnote}{-1}%
	\endgroup
}
\spanishdatedel
\decimalpoint

\def\LOGOUNI{%
	\begin{picture}(0,0)\unitlength=1cm
	\put (-13.5,-2.5) {\includegraphics[width=2.8cm]{logouni}}
	\end{picture}
}

\def\LOGOCTIC{%
	\begin{picture}(0,0)\unitlength=1cm
	\put (-8,-1.2) {\includegraphics[height=1cm]{logocticblack}}
	\end{picture}
}

\newcommand{\mychar}{%
	\begingroup\normalfont
	\includegraphics[height=\fontcharht\font`\B]{Octocat.png}%
	\endgroup
}

\begin{document}

\begin{center}
	\sffamily\bfseries\scshape
	{\Large UNIVERSIDAD NACIONAL DE INGENIERÍA}\LOGOUNI\\
	Centro de Tecnologías de la Información y Comunicaciones\LOGOCTIC\\
\end{center}

\vspace{.8cm}

\begin{center}\sffamily\bfseries\large
	Programación en \texttt{C++} \\
	TEMA $\bm{7}$: Arreglo de caracteres.
\end{center}

\vspace{.5cm}
\noindent
\makebox[\textwidth]{Nombres y apellidos:\enspace\hrulefill}
\vspace{0.2in}
\makebox[\textwidth]{Nombres y apellidos del instructor:\enspace MSc. Víctor Melchor Espinoza.\hfill}

\section*{Ejercicios obligatorios}

\begin{questions}

\question Elabore un programa que solicite una contraseña para el acceso a un sistema. Cada vez que el usuario no adivine se mostrará el mensaje. \textcolor{green}{!Contraseña incorrecta¡}. El usuario tiene tres oportunidades para adivinar la clave. Si adivina la clave se mostrará el mensaje \textcolor{green}{Bienvenido al sistema} y se terminará el programa. En el caso que tenga tres intentos fallidos, entonces debe mostrar el mensaje \textcolor{green}{Cuenta bloqueada} y a continuación finaliza el programa.

\begin{solution}
	\begin{listing}[H]
		\footnotesize
		\inputminted{cpp}{1.cc}
		\caption{Programa \texttt{1.cc}.}
		\label{lst:7.1}
\end{listing}
\end{solution}
% TODO 5
\question Un grupo de inteligencia militar desea codificar los mensajes secretos de tal forma que no puedan ser interpretados con una lectura directa; para lo cual han establecido las siguientes reglas:

\begin{parts}
	\part Todo mensaje debe tener sus letras en mayúsculas.
	\part Reemplazar cada letra por la que sigue según el abecedario excepto la 'Z' que deberá reemplazar con la letra 'A'.
	\part Reemplazar cada dígito encontrado por el que le precede según su orden excepto el '0'que deberá reemplazarse por el '9'.
\end{part}
Ingresar un mensaje y mostrar su codificación secreta resultante.

\begin{solution}
	\begin{listing}[H]
		\footnotesize
		\inputminted{cpp}{5.cc}
		\caption{Programa \texttt{5.cc}.}
		\label{lst:7.2}
\end{listing}
\end{solution}

\end{questions}

\begin{flushright}\bfseries
Centro de Tecnologías de la Información y Comunicaciones (CTIC)\\[2mm]
\today\unmarkedfntext{Código disponible en \href{https://github.com/carlosal1015/Cpp-Programming}{\mychar{}}.}
\end{flushright}

\end{document}